% ==========================================
% N5 - DIEU KHIEN DUONG DI OSPF BANG COST
% Nhom K = 7126
% ==========================================

\section{N5 -- Điều khiển đường đi OSPF bằng cost}
\label{sec:N5}

\subsection{Mục tiêu}

Nhiệm vụ N5 điều chỉnh cost OSPF để lưu lượng từ B1
đến B3 đi qua B2 thay vì sử dụng trực tiếp liên kết L8.
Nhóm đối chiếu traceroute trước và sau thay đổi,
đồng thời phân tích ảnh hưởng của đường đi nội bộ
đến lưu lượng giữa các AS.

Thử nghiệm sử dụng mô hình OSPF area 0 trong AS B
đã triển khai ở Phần C. Hai đường giữa B1 và B3 gồm
đường trực tiếp trên L8 và đường qua B2 trên L6, L7.
Đích kiểm tra là Loopback0 của B3, có địa chỉ
\texttt{10.171.26.6/32}.

\subsection{Triển khai cấu hình}

\subsubsection{Nguyên tắc lựa chọn đường đi}

OSPF lựa chọn đường có tổng cost nhỏ hơn.
Do đó, có thể tăng cost của cổng đi ra L8 trên B1
để ưu tiên đường B1 $\rightarrow$ B2 $\rightarrow$ B3.
Theo chú thích đi kèm ảnh trong tài liệu minh chứng,
cost được thay đổi từ \textbf{1 lên 100}.

Gọi $C_{\text{L6}}$ là cost từ B1 tới B2 trên L6,
$C_{\text{L7}}$ là cost từ B2 tới B3 trên L7 và
$C_{\text{L8}}$ là cost từ B1 tới B3 trên L8.
Điều kiện để ưu tiên đường qua B2 là:

\begin{equation}
C_{\text{L6}} + C_{\text{L7}} < C_{\text{L8}} = 100.
\end{equation}

Khi xét đường tới Loopback0 của B3,
hai phương án còn có cùng thành phần cost tới interface đích,
nên thành phần này không làm thay đổi phép so sánh.
Hai ảnh traceroute không hiển thị cost riêng của L6 và L7;
kết quả được đánh giá qua đường đi thực tế quan sát được.

\subsubsection{Cấu hình trên B1}

Theo bảng quy hoạch địa chỉ, L8 nối FastEthernet1/0 của B1
với FastEthernet0/1 của B3.
Để đặt cost L8 phía B1 bằng 100, sử dụng cấu hình:

\begin{lstlisting}
B1-7126#configure terminal
B1-7126(config)#interface FastEthernet1/0
B1-7126(config-if)#ip ospf cost 100
B1-7126(config-if)#end
\end{lstlisting}

Lệnh \texttt{ip ospf cost} đặt chi phí gửi gói qua interface.\footnote{Cisco,
\href{https://www.cisco.com/c/en/us/td/docs/ios-xml/ios/iproute_ospf/command/iro-cr-book/ospf-a1.html}
{Cisco IOS IP Routing: OSPF Command Reference -- ip ospf cost}.}
Việc tăng cost cho phép thay đổi mức ưu tiên đường đi
trong khi vẫn duy trì liên kết L8.
Sau khi OSPF cập nhật đường đi, chạy lại cùng lệnh traceroute
trên B1 để so sánh với trạng thái ban đầu.

\subsection{Kiểm chứng kết quả}

\subsubsection{Trước khi thay đổi cost}

Trên B1, nhóm kiểm tra đường tới Loopback0 của B3:

\begin{lstlisting}
B1-7126#traceroute 10.171.26.6
\end{lstlisting}

\partzeroimage
{images/advanced/N5/B1_traceroute_before.png}
{Traceroute từ B1 tới Loopback0 của B3 trước khi thay đổi cost}
{fig:N5-trace-before}

Hình~\ref{fig:N5-trace-before} ghi nhận một hop
\texttt{10.71.26.30}, là địa chỉ cổng Fa0/1 của B3
trên L8. Như vậy, đường đi ban đầu là
\textbf{B1 $\rightarrow$ B3}, sử dụng liên kết trực tiếp L8.
Traceroute tới địa chỉ Loopback0 có thể nhận phản hồi
từ địa chỉ interface của router đích; ở đây,
\texttt{10.71.26.30} vẫn thuộc B3.

\subsubsection{Sau khi tăng cost lên 100}

Nhóm thực hiện lại lệnh \texttt{traceroute 10.171.26.6}
trên cùng router B1 sau khi thay đổi cost.

\partzeroimage
{images/advanced/N5/B1_traceroute_after.png}
{Traceroute từ B1 tới B3 sau khi tăng cost L8 phía B1 lên 100}
{fig:N5-trace-after}

Hình~\ref{fig:N5-trace-after} ghi nhận hai hop:
\begin{itemize}
    \item Hop 1: \texttt{10.71.26.22}, địa chỉ Fa0/0
    của B2 trên L6.
    \item Hop 2: \texttt{10.71.26.26}, địa chỉ Fa0/0
    của B3 trên L7.
\end{itemize}

Kết quả xác nhận đường đi đã chuyển sang
\textbf{B1 $\rightarrow$ B2 $\rightarrow$ B3},
sử dụng L6 và L7, đáp ứng yêu cầu N5.

\begin{table}[H]
\centering
\caption{So sánh đường đi từ B1 tới B3 trong nhiệm vụ N5}
\label{tab:N5-comparison}
\small
\renewcommand{\arraystretch}{1.3}
\begin{tabularx}{\textwidth}{l X X}
\toprule
\textbf{Nội dung} & \textbf{Trước thay đổi} & \textbf{Sau thay đổi} \\
\midrule
Cost L8 phía B1 & 1 & 100 \\
Đường đi & B1 $\rightarrow$ B3 & B1 $\rightarrow$ B2 $\rightarrow$ B3 \\
Liên kết sử dụng & L8 & L6, L7 \\
Hop quan sát & \texttt{10.71.26.30} &
\texttt{10.71.26.22}, \texttt{10.71.26.26} \\
\bottomrule
\end{tabularx}
\end{table}

\subsection{Nhận xét và ảnh hưởng tới lưu lượng giữa các AS}

Kết quả cho thấy OSPF có thể chọn đường đi qua thêm một router
khi đường đó có tổng cost thấp hơn.
Các số đo mili giây trong traceroute là thời gian phản hồi
của từng lần thăm dò, không phải giá trị cost OSPF.

Trong mô hình, B1 là router biên nối AS A qua L4,
còn B3 nối AS C qua L9. Khi lưu lượng từ AS A tới AS C
được BGP chọn đi xuyên AS B qua B1 và B3,
việc điều chỉnh OSPF có thể thay đổi đoạn đường nội bộ:

\begin{center}
Trước thay đổi: A2 $\rightarrow$ B1 $\rightarrow$ B3 $\rightarrow$ C1.

Sau thay đổi: A2 $\rightarrow$ B1 $\rightarrow$ B2
$\rightarrow$ B3 $\rightarrow$ C1.
\end{center}

Với cấu hình \texttt{next-hop-self} ở Phần D,
đường OSPF tới Loopback0 của router biên được dùng
để phân giải next-hop iBGP. Vì vậy, điều chỉnh cost
có thể chuyển lưu lượng từ L8 sang L6, L7 trong khi
AS\_PATH vẫn đi qua cùng các AS.
Đây là phân tích từ topology và cấu hình BGP;
hai ảnh minh chứng trực tiếp xác nhận đường B1 tới B3,
chưa có traceroute giữa hai LAN để kiểm chứng toàn tuyến.

Cost được đặt trên interface theo chiều gửi.
Nếu chỉ tăng cost phía B1, chưa thể kết luận
chiều B3 về B1 cũng đi qua B2; chiều về có thể vẫn dùng L8.
Việc kiểm chứng N5 vì vậy tập trung vào chiều
B1 tới B3 đúng yêu cầu đề bài.

\newpage
